\chapter{Semi-automatic Trader}

I N THIS FINAL PART OF THE BOOK I'LL SHOW YOU HOW TO PUT THE framework into action for three specific examples. This first chapter is for semiautomatic traders, who make their own discretionary \textbf{forecasts} about price movements, but then use my systematic \textbf{framework} to manage their capital and position risk.

\section*{Chapter overview} \phantomsection \addcontentsline{toc}{section}{Chapter overview}

\begin{chapteroverviewtable}
\chapteroverviewitem{Who are you?}{Introducing the semi-automatic trader.}
\chapteroverviewitem{Using the framework}{How you will use the systematic framework as a semi-automatic trader.}
\chapteroverviewitem{Process}{The process you need to follow each morning and as you are trading throughout the day.}
\chapteroverviewitem{Trading diary}{A diary showing how you could have traded the markets as a semi-automatic trader in late 2014.}
\end{chapteroverviewtable}

\section*{Who are you?} \phantomsection \addcontentsline{toc}{section}{Who are you?}

As a semi-automatic trader you want to make your own calls on the market, but within a systematic \textbf{framework}. With stop losses, risk targeting and position sizing taken care of you can focus on getting the buy or sell decision correct. This will give you the best of both worlds -- your human ability to interpret and process information, combined with a system giving the correct amount of risk.

It will not be easy sticking to the framework. The system may force you to trade, or prevent you from trading, when you would rather do otherwise. However there are significant benefits in sticking to a consistent strategy. I strongly recommend that you design a system you're comfortable with and then do not deviate from it.

In the specific example discussed in this chapter we'll assume you're going to be trading spread bets on equity indices and FX, as reasonable sized amateur investors with a notional trading capital of £100,000. You could be trading full-time, or on a part-time basis for perhaps an hour each morning. However you can apply this framework to any market where leverage is possible and to any account size.

As a semi-automatic trader you'll be trading sporadically as opportunities arise. This presents some challenges in using my systematic framework, which normally requires a series of calculations to be performed each time you trade.

To reduce the workload you'll first do some daily housekeeping on your current portfolio: closing out any positions that have hit stop losses, adjusting estimates of \textbf{trading capital} and \textbf{price volatility}, and doing any trades on existing positions that are now required. You'll then put on any new trades which look promising based on overnight moves.

If you're trading part-time you will leave stop loss orders with your broker and go to your day job. However if you're trading full-time you'll continue to watch the markets, potentially putting further trades on, and closing positions that hit stop losses. You won't however repeat the more time-consuming housekeeping tasks until the next morning.

\section*{Using the framework} \phantomsection \addcontentsline{toc}{section}{Using the framework}

\subsection*{Instrument choice: size and costs}

As \textbf{semi-automatic traders} you don't trade a fixed set of instruments. At any given time you will have a small number of positions, or ‘bets', open. Those positions are drawn from a larger pool of \textbf{instruments} that you've formed opinions on.

You'll be using quarterly spread bets, with all bets in multiples of £1 per point. For example if you go long the FTSE at £2 per point, and it goes from 6500 to 6501, then you'd earn £2.\footnote{There are several other important aspects of spread betting that I don't have space to discuss here. I assume that if you intend to actually trade the system shown here you are already familiar with spread betting. If not you should consult the books I've recommended in appendix A, or similar resources.} The bet size will be determined by the broker's minimum bet size, usually between £1 and £10 depending on the instrument. In line with the advice in chapter twelve, ‘Speed and Size', (page 200) you should check that any position you're putting on would be at least four instrument blocks with a maximum forecast of 20.

You should also check that the \textbf{standardised costs} of trading any potential instrument will be 0.01 \textbf{Sharpe ratio} units or less, for reasons explained in detail below. Finally you should avoid very low volatility instruments, for the numerous reasons enumerated throughout part three.

\subsection*{Forecasts and stop losses}

As I discussed in chapter seven, ‘Forecasts', (page 115) you'll need to have a graduated opinion on likely price movements whenever you decide to open a new position or ‘bet'. These need to be translated into quantifiable \textbf{forecasts}, as in table 37. Note if you weren't using spread bets or other derivatives, and couldn't \textbf{short sell} assets, you'd need to limit yourself to making positive forecasts.

\rawtablecaption{TRANSLATING OPINIONS INTO A QUANTIFIED FORECAST\protect\footnotemark}
\footnotetext{This is identical to table 16 and is repeated here for ease of reference.}
\begingroup \small \setlength{\tabcolsep}{2pt} \renewcommand{\arraystretch}{1.15} \begin{tabularx}{\textwidth}{@{}*{9}{>{\centering\arraybackslash} X}@{}}
\toprule
\shortstack{Very\\strong\\sell} & \shortstack{Strong\\sell} & Sell & \shortstack{Weak\\sell} & Neutral & \shortstack{Weak\\buy} & Buy & \shortstack{Strong\\buy} & \shortstack{Very\\strong\\buy} \\
\midrule
-20 & -15 & -10 & -5 & 0 & 5 & 10 & 15 & 20 \\
\bottomrule
\end{tabularx}
\endgroup

I strongly recommend that you do not change your forecast once a bet is open. Otherwise you might be tempted to take profits too early, or double up on losses. You can sometimes add new bets on top of an existing position, of which more later. But you will be using a systematic trailing stop loss rule to close all your positions; no other exit rules are permitted.

The stop loss rule is loosely related to the ‘early loss taker’ system I’ve used for examples in previous chapters. It uses stops set at a multiple of the current value of the daily \textbf{standard deviation} of prices.

\noindent{\small \setlength{\tabcolsep}{4pt} \renewcommand{\arraystretch}{1.15} \begin{tabularx}{\textwidth}{@{} p{0.24\textwidth} Y@{}}\textbf{Parameter \(X\)} & You should set a parameter \(X\) which will be related to how long you want to hold positions for and your trading costs. I recommend \(X = 4\), for reasons I explain below. \\
\textbf{Volatility units} & The expected standard deviation of daily price changes, in price points. Note we use volatility in price points, not percentage points as normal. The volatility in price points is equal to the percentage point volatility (price volatility as defined in chapter ten, ‘Position sizing’, on “Price volatility” on page 155), multiplied by the current price. In this chapter you’ll find this by eyeballing a one month chart of recent price changes (as discussed in chapter ten, page 155), since you’re probably going to be using charts to decide on your entries. You need to update this estimate throughout the life of the trade and adjust the stop accordingly. As suggested in chapter twelve, you don’t need to update your volatility estimate unless it has changed by more than 25\% from the previous value. For example as I write this the standard deviation of oil prices is \(\$1.5\) per day. \\
\textbf{Trailing stop loss when long} & If the price of the instrument falls by more than \(X\) volatility units from the high achieved since entry then you should close your position. As an example suppose the high of crude was \(\$70\), so with \(X = 4\) and volatility units of \(\$1.5\), that implies a stop loss of \(\$70 - (4 \times \$1.5) = \$64\). \\
\textbf{Trailing stop loss when short} & You will generate a closing trade when a price rises by more than \(X\) volatility units from the low reached after entry. So again for crude if the recent low was \(\$30\), then that implies a stop loss of \(\$30 + (4 \times \$1.5) = \$46\). \\
\end{tabularx}
}

\smallskip \noindent\textit{* \(X\) is equivalent to the \(B\) parameter in the generic ‘A and B’ system which I describe in appendix B.}

What value of \(X\) should you use? It depends on whether you have a short or long-term forecasting horizon, and on the costs of your instruments. So with a shorter forecasting horizon you’ll cut more quickly, hold positions for less time, and it will be more expensive to trade. It’s important to match the forecasting horizon and the tightness of your stop loss. There is no point making a forecast of prices for the next six months if you’re likely to be stopped out by next Tuesday. Table 38 shows the average holding period for a given value of \(X\).\footnote{I produced this table using both artificial and real market data drawn from a number of asset classes. Because we obviously can’t simulate a discretionary rule I used a random entry rule. Interestingly this is actually profitable in real data, because even with a random entry we can still capitalise on trends by using stop losses.}

\rawtablecaption{HOW DOES THE X PARAMETER IN YOUR STOP LOSS RULE DETERMINE YOUR AVERAGE HOLDING PERIOD? SMALLER VALUES OF X MEAN TIGHTER STOP LOSSES AND SHORTER HORIZONS, BUT REQUIRE CHEAPER INSTRUMENTS} \noindent{\small \setlength{\tabcolsep}{4pt} \renewcommand{\arraystretch}{1.15} \begin{tabularx}{\textwidth}{@{} p{0.18\textwidth} p{0.24\textwidth} p{0.28\textwidth} Y@{}}
\toprule
\textbf{Value of \(X\)} & \textbf{Average holding period} & \textbf{Turnover, round trips per year} & \textbf{Maximum cost SR} \\
\midrule
1 & 4 days & 64 & 0.0013 \\
2 & 9 days & 29 & 0.0027 \\
3 & 17 days & 15 & 0.0053 \\
4 & 6.5 weeks & 8 & 0.01 \\
8 & 13 weeks & 4.0 & 0.02 \\
10 & 26 weeks & 2.0 & 0.04 \\
\bottomrule
\end{tabularx}
}

The table shows for each value of X (rows): the average holding period, in business days or weeks; the implied turnover in round trips per year; and the maximum instrument cost if we pay no more than 0.08 Sharpe ratio units per year. The turnover includes the additional trading we get from changes in price volatility and also assumes we don't trade position changes of less than 10\% (position inertia).

As you saw in chapter twelve, ‘Speed and Size', you need to adjust X given the cost of the instruments you're trading. The final column of table 38 shows a suggested maximum feasible \textbf{standardised cost}, in \textbf{Sharpe ratio (SR)} units per year, for a given value of X. This is calculated using my recommended \textbf{speed limit} of spending 0.08 SR units per year on costs, from chapter twelve.

Also in chapter twelve (page 183) I suggested you use a costs figure of 0.01 SR units as the standardised cost for all of your spread bets. The table shows you could safely use an X of 4 or higher with spread bets, although if you're trading cheap futures contracts you could go quicker.

I don't advise using different stop loss systems for different instruments as this complicates life and can lead to mistakes. So the value of X should be set conservatively using the cost of the most expensive instrument you're likely to trade.

In this chapter we are going to use an X of 4, which means I assume you're predicting trends that last for just over six weeks.

Note that unlike in the original ‘early loss taker' system you won't be using a profit target. Why not? My own research shows no evidence that systematic profit targets work consistently.\footnote{There are some situations where closing trades before a stop loss is hit makes sense, such as \textbf{relative value} or \textbf{mean reverting} strategies, but in my opinion it is better to build the closing rule into a fully systematic trading rule.} This is because most markets exhibit trending behaviour, and a profit target will get you out of trends too early.

To summarise you mustn't close your trade until you hit a stop loss; then you must close it. Sticking to this rule will be difficult, but will mean that your risk is controlled and that your returns will have the favourable positive skew of \textbf{trend following} traders.

\subsection*{Volatility targeting}

I've already set the initial \textbf{trading capital} of this example at £100,000. But how aggressive should you be with this money -- what should your \textbf{percentage} \textbf{volatility target} be? You need to refer back to chapter nine, ‘Volatility Targeting'. In this example I am going to assume that your \textbf{Sharpe ratio} (SR) will be at least 0.30 after costs, which I calculate below as 0.08 SR units. This is comfortably under my recommended maximum expectation of 0.50 (from page 46). Referring to column D of table 26 (page 148) an SR of 0.30 gives a risk percentage of 15\%, and I'm assuming you can cope with the pain that implies.
This implies you'll have an annual \textbf{volatility target} of £100,000 × 15\% = £15,000 and a daily target of £15,000 ÷ 16 = £937.50.\footnote{As before to go from annual to daily volatility you need to divide by the ‘square root of time’, which assuming 256 business days in a year is 16.} Because you're using spread bets there is no concern about hitting a relatively modest risk percentage of 15\%. As always you should avoid low volatility assets where the required leverage will be too high.

You'll be measuring your account value daily, recalculating \textbf{trading capital}, and adjusting your \textbf{volatility target} accordingly. This might trigger trades on existing positions and the new volatility target will be used to size any new bets you make.

\subsection*{Position sizing}

You can now think about ‘Position sizing', as discussed in chapter ten. First you need to measure the recent percentage \textbf{standard deviation} of returns (\textbf{price volatility}) in each market to determine the expected daily variation in the value of one \textbf{instrument block} (the instrument currency volatility).

I'd suggest you use the eyeball technique described in chapter ten (page 155), which you're already using to calculate the volatility estimate needed for stop loss levels. This should be done daily on any instruments where you have positions, and ideally on anything you might place bets on today.\footnote{Working out the instrument currency volatility for a potential trade in advance makes the process more efficient but might not always be possible, in which case you will need to do your estimation just before you trade.} Changes in price volatility for your existing positions might trigger trades.

From table 38 (page 212) with the suggested X of 4 you would expect to get a turnover of 8 after accounting for the effects of sizing your position. With \textbf{standardised costs} of 0.01 SR units for spread bets this gives us an annual cost of 8 × 0.01 = 0.08 \textbf{Sharpe ratio} (SR) units, which just squeaks in at the maximum level of 0.08 SR that I recommended in chapter twelve, ‘Speed and Size'.

With 15\% annualised volatility that equates to a performance drag from costs of 0.08 × 15\% = 1.2\% a year. This is more than most \textbf{passive ETFs}, but significantly cheaper than a hedge fund. You'd also have to pay perhaps 0.6\% for rolling your quarterly bets.\footnote{Assuming the full spread is charged on rolls.}

You can save on costs by following the advice in the chapter twelve, ‘Speed and Size', and not adjusting your estimate of \textbf{instrument currency volatility} if the previous level looks about right, say within 25\% of the current value.

As you're spread betting, the value of a 1\% move (the \textbf{block value}) depends on the bet size in pounds per point, so you won't need to use exchange rates even when betting on non-UK markets. So for example with crude oil at \$65, a 1\% move of \$0.65 or 0.65 points, at a broker's minimum of £10 a point, will have a value of £6.50.

With this information you can use your \textbf{volatility target} to work out the \textbf{volatility scalars} every morning to use for both existing and potential trades.

\subsection*{Portfolios}

To size your bets you need to determine your preferred maximum and average number of concurrent bets, as I explained in chapter eleven, ‘Portfolios'. These terms are self explanatory: the maximum number of bets is the most positions you will have on at any given time, and the average is what you expect to have on a typical day. As I said in chapter eleven, the \textbf{instrument weight} you use to size each of your positions will be 100\% divided by the maximum number of bets.

The \textbf{instrument diversification multiplier} will be the maximum number of bets divided by the average. This multiplier ensures the total average absolute forecast will be 10 with the average number of bets on, which ensures your forecasting is consistent with the rest of the systematic \textbf{framework}. I recommended that the multiplier shouldn't be greater than 2.5, which limits the value of the maximum relative to the average.

If you already have your maximum number of bets on, and you desperately need to make a new bet, then sorry: you can't. Wait until one of your existing bets has stopped out. If they all continue to be profitable then perhaps your new bet wasn't necessary after all. Although this may seem rigid, sticking to this ensures that you have the correct risk on for each bet and it will stop you from putting on dangerous numbers of simultaneous bets. Make sure you set your maximum high enough when designing your system to cover your likely appetite for extra bets.

Earlier I said you shouldn't change the forecast on an existing bet. However you can add to positions by placing a new, separate, bet on an instrument which you're already holding. This allows pyramiding positions, such as buying into a strengthening trend. The new bet is completely distinct, will have its own stop loss, and could be closed at a different time. A new bet on an existing position counts towards your maximum number of bets.

This can't be used as a loophole to close positions; you can't put on a new bet which will make your existing position smaller or reverse it. Finally, to avoid concentrating your risk too heavily I strongly recommend that you don't have a total of more than 40 forecast units bet on any asset at once, added up over all the bets on that instrument. So a forecast of +10 on your first bet, followed by +25 on your second for the same instrument would be acceptable. But you couldn't place a third bet of more than +5.

In this chapter I am going to assume you have a maximum of four positions, with an average of three. So each position will have an instrument weight of 100\% ÷ 4 = 25\%, and the diversification multiplier is 4 ÷ 3 = 1.33.

As you saw in chapter nine, ‘Volatility targeting', (page 150) the use of systematic stop losses allows you to compare my framework with traditional money management systems, where you put a fixed percentage of your capital at risk. Each bet you make puts at risk a percentage of capital equal to: (

X × forecast × percentage volatility target) ÷ (10 × 16 × average number of bets)

In this example with an X of 4, the average forecast of 10, a 15\% volatility target and an average of three positions you are putting (4 × 10 × 15\%) ÷ (10 × 16 × 3) = 1.25\% of your capital at risk on average per bet. With the maximum forecast of 20 you'd be risking twice that, 2.5\%. This is relatively large, but remember that you expect to be holding these positions for around six weeks.

\subsection*{Trading in practice}

You should first check to see if any stop losses on existing positions have been hit, and trade on those immediately if you didn't leave stop losses with your broker.\footnote{I'm mostly agnostic about whether stop loss orders should be left with brokers, or entered only when prices hit the relevant level. However, unless you are watching the markets all day it is safer to leave stop loss orders in place. Pre-entered orders also prevent you from ignoring an exit and staying in the trade.} Then you can look at potential new entries and come up with a numeric forecast for them. On every new trade you'll need to calculate a stop loss based on the value of X, the \textbf{price volatility} and the entry price.

Stop losses on existing positions will also need to be adjusted if price volatility changes. They also need to be moved up if prices reach new highs (for longs), and down for new lows (on shorts).

Part-time daily traders should then update the stop loss limits they have left with their brokers and go to work at their real jobs. If you're still trading intra-day you will continue to monitor the market for new entries and to see if stops have been hit. You could also update stop loss levels intra-day if significant new levels are reached, but with the relatively large value of X I've suggested for this chapter this probably isn't necessary.

\section*{Process} \phantomsection \addcontentsline{toc}{section}{Process}

Here is the process to follow when running this system. The trading diary in the next section illustrates in more detail how the calculations are done.

\subsection*{Daily housekeeping}

\noindent{\small \setlength{\tabcolsep}{2pt} \renewcommand{\arraystretch}{1.15} \begin{tabularx}{\textwidth}{@{} p{0.29\textwidth} Y@{}}\textbf{Check stops} & If not using stop orders left with brokers, check to see if any stop levels have been hit on existing positions, and if necessary do closing trades. \\
\textbf{Get account value} & Get today's account value and calculate your current trading capital. Your annualised cash volatility target is current capital multiplied by percentage volatility target. In this chapter we use 15\%. \\
\textbf{Daily cash volatility target} & Daily cash volatility target equals one-sixteenth of annualised cash volatility target. As usual this is the ‘square root of time' rule; with around 256 business days in a year you should divide by 16 to go from annual to daily risk. \\
\textbf{Get latest prices} & Get charts for all instruments you have positions on and ideally anything you might wish to trade today. \\
\textbf{Calculate daily price volatility} & Using a one month chart eyeball the price volatility of each instrument you own, or might trade today. With FX rate of 1 and relevant block values convert this to instrument value volatility. For existing positions if the volatility level you used yesterday is not within 25\% of today's level then update your estimate; otherwise stick with yesterday's estimate. For potential positions you should always use the most up-to-date estimate. \\
\textbf{Recalculate stops} & Update stops on existing instruments so they trail new highs for longs or new lows for shorts, and account for the latest value of price volatility. \\
\textbf{Volatility scalar} & This is equal to the daily cash volatility target divided by instrument value volatility for each instrument. \\
\textbf{Desired subsystem position on an existing bet} & Original forecast multiplied by volatility scalar, divided by 10. \\
\textbf{Desired portfolio instrument position on an existing bet} & Desired subsystem position multiplied by instrument weight and by instrument diversification multiplier. With the maximum of four bets used in the example, and an average of 3 bets, you should use an instrument weight of 25\% and a multiplier of 1.33. \\
\textbf{Desired target position on an existing bet} & Portfolio instrument position rounded to the nearest whole block. \\
\textbf{Issue trades on existing positions} & Compare current position to target position. If out by more than 10\% then issue adjusting orders, using position inertia. \\
\textbf{Check setups} & Check charts and news for potential new bets. If any exist go to instructions for new position opening. \\
\end{tabularx}
}

\subsection*{Further intra-day process (for full-time traders only)}

\noindent{\small \setlength{\tabcolsep}{2pt} \renewcommand{\arraystretch}{1.15} \begin{tabularx}{\textwidth}{@{} p{0.29\textwidth} Y@{}}\textbf{Check stop losses} & Check to see if any stop losses have been triggered. If not using broker stops then issue closing trades. \\
\textbf{Check setups} & Check charts and news for potential trade setups. If any exist go to instructions for new position opening. \\
\end{tabularx}
}

\subsection*{New position opening}

\noindent{\small \setlength{\tabcolsep}{2pt} \renewcommand{\arraystretch}{1.15} \begin{tabularx}{\textwidth}{@{} p{0.29\textwidth} Y@{}}\textbf{Calculate forecast} & Use table 37 (page 211) to calculate your forecast for the instrument. \\
\textbf{Is it allowed?} & Do you already have the maximum number of bets on? For the example system this is four. Is the bet on an instrument with an existing position? If so it must not reduce or reverse the position, or put the total absolute forecast across all bets for this instrument over 40. \\
\textbf{Calculate instrument value volatility} & If not pre-calculated in the morning, using a one month chart eyeball the price volatility of each instrument. With FX rate of 1 and relevant block values convert this to instrument value volatility. \\
\textbf{Calculate stop loss} & Use \(X\) multiplied by the daily price volatility in price points to find where to set the stop loss relative to the entry price. In this chapter I've used an \(X\) of 4. \\
\textbf{Volatility scalar} & This morning's daily cash volatility target divided by instrument value volatility for the instrument. \\
\textbf{Subsystem position} & Forecast multiplied by volatility scalar, divided by 10. \\
\textbf{Portfolio position} & Subsystem position multiplied by instrument weight, 25\% in the example, and then multiplied by instrument diversification multiplier, 1.33 in the example. \\
\textbf{Rounded target position} & Portfolio instrument position rounded to nearest block. \\
\textbf{Trade} & You should now put on the trade and once completed calculate the initial stop loss relative to the entry price. With judicious use of spreadsheets the above steps can be completed in a few moments. \\
\end{tabularx}
}

\section*{Trading diary} \phantomsection \addcontentsline{toc}{section}{Trading diary}

Here is some paper trading I did using the semi-automatic trading system. All the calculations here have been done with a spreadsheet, which is available from my website. Prices and other values are rounded to make the example clearer.

\subsection*{15 October 2014}

Typical minimum spread bets are £10 a point on crude, £5 on US S\&P 500 equities and £1 on the European Euro Stoxx equity index. These and the price determine the \textbf{block value.} Note that the block values are all in British pounds, the same currency as my account, even though the prices are quoted in other currencies, as we're betting in £1 per point. So the FX rates are all 1. Starting \textbf{trading capital} is £100,000; annual \textbf{volatility target} at 15\% is £15,000 and starting daily volatility target is one sixteenth of that, £937.50.

\noindent{\scriptsize \setlength{\tabcolsep}{2pt} \renewcommand{\arraystretch}{1.12} \begin{tabularx}{\textwidth}{@{} p{0.36\textwidth}*{3}{R}@{}}
\toprule
& Crude & S\&P 500 & Euro Stoxx \\
\midrule
Daily volatility target (A) & £937.50 & £937.50 & £937.50 \\
Price (B) & 83 & 1880 & 2900 \\
Bet size, per point (C) & £10 & £5 & £1 \\
FX (D) & 1 & 1 & 1 \\
Block value (E) & £8.30 & £94 & £29 \\
\emph{C × B ÷ 100} & & & \\
Price volatility, points (F) & 1 & 20 & 50 \\
Price volatility, \% (G) & 1.20 & 1.06 & 1.72 \\
\emph{F ÷ B} & & & \\
Instrument currency volatility (H) & £10 & £100 & £50 \\
\emph{G × E} & & & \\
Instrument value volatility (I) & £10 & £100 & £50 \\
\emph{H × D} & & & \\
Volatility scalar (J) & 93.75 & 9.375 & 18.75 \\
\emph{A ÷ I} & & & \\
\bottomrule
\end{tabularx}
}

For my forecasts I'm feeling very bullish on US equities, bullish on crude and bearish on European equities.

\noindent{\scriptsize \setlength{\tabcolsep}{2pt} \renewcommand{\arraystretch}{1.12} \begin{tabularx}{\textwidth}{@{} p{0.36\textwidth}*{3}{R}@{}}
\toprule
& Crude & S\&P 500 & Euro Stoxx \\
\midrule
Forecast (K) & 10 & 15 & -10 \\
Subsystem position, blocks (L) & 93.75 & 14.06 & -18.75 \\
\emph{K × J ÷ 10} & & & \\
Instrument weight (M) & 25\% & 25\% & 25\% \\
Instrument diversification multiplier (N) & 1.33 & 1.33 & 1.33 \\
Portfolio instrument position, blocks (O) & 31.25 & 4.687 & -6.2 \\
\emph{L × M × N} & & & \\
Rounded target position, blocks (P = round O) & 31 & 5 & -6 \\
Entry price (Q) & 83 & 1880 & 2900 \\
Stop loss offset (R) & 4 & 80 & 200 \\
\emph{F × X with X = 4} & & & \\
Stop loss (S) & 79 & 1800 & 3100 \\
\emph{Q plus or minus R} & & & \\
\bottomrule
\end{tabularx}
}

So I go long crude at 31 × £10 = £310 a point with a \$79 stop loss, long S\&P 500 at 5 × £5 = £25 per point with an 1800 stop loss, and short Euro Stoxx at 6 × £1 = £6 per point with a 3100 stop loss.

\subsection*{29 October 2014}

\noindent{\scriptsize \setlength{\tabcolsep}{2pt} \renewcommand{\arraystretch}{1.12} \begin{tabularx}{\textwidth}{@{} p{0.36\textwidth}*{3}{R}@{}}
\toprule
& Crude & S\&P 500 & Euro Stoxx \\
\midrule
Forecast & 10 & 15 & -10 \\
Entry price & 83 & 1880 & 2900 \\
Current price & 81 & 1980 & 3020 \\
High (Low) since entry (T) & 83 & 1980 & (2900) \\
Stop loss (S = T plus or minus R) & 79 & 1900 & 3100 \\
Current position, blocks & 31 & 5 & -6 \\
Gain (Loss) & (£620) & £2,500 & (£720) \\
Portfolio instrument position, blocks (O) & 31.5 & 4.73 & -6.3 \\
Rounded target position, blocks (P) & 32 & 5 & -6 \\
\bottomrule
\end{tabularx}
}

No stops triggered and the trailing stop loss for S\&P 500 has been moved up. The net profit in my account is £1,160, so current capital is £101,160 and the daily volatility target is £948. The \textbf{price volatility} and \textbf{instrument value volatility} hasn't changed on any instrument by more than 25\%, so I haven't adjusted these. Current positions are still within 10\% of the target; there is no need to buy one more \textbf{instrument block} of crude to get to a rounded position of 32.

I feel very strongly about the S\&P bet now and I'm going to add to my position. So I will bet a separate +25 forecast, taking me up to the maximum of +40 for all S\&P 500 bets. It will also help hedge my Euro Stoxx short, which is doing very badly though hasn't yet hit it's stop. I now have four bets on, two of which are in S\&P 500, so I can't place any more.

\noindent{\scriptsize \setlength{\tabcolsep}{2pt} \renewcommand{\arraystretch}{1.12} \begin{tabularx}{\textwidth}{@{} p{0.40\textwidth}*{2}{R}@{}}
\toprule
& S\&P 500 (2) & S\&P 500 (2) \\
\midrule
Daily volatility target (A) & £948 & Forecast (K): 25 \\
Price (B) & 1980 & Subsystem position (L): 23.7 \\
\emph{K × J ÷ 10} & & \\
Bet size, per point (C) & £5 & Instrument weight (M): 25\% \\
FX (D) & 1 & Instrument diversification multiplier (N): 1.33 \\
Block value (E) & £99 & Portfolio instrument position (O): 7.88 \\
\emph{C × B ÷ 100} & & \emph{L × M × N} \\
Price volatility, points (F) & 20 & Rounded target position, and trade (P = round O): 8 \\
Price volatility, \% (G) & 1.01 & Entry price (Q): 1980 \\
\emph{F ÷ B} & & \\
Instrument currency volatility (H) & £100 & Stop loss offset (R): 80 \\
\emph{F × X with X = 4} & & \\
Instrument value volatility (I) & £100 & Stop loss (S): 1900 \\
\emph{G × E} & & \emph{Q plus or minus R} \\
Volatility scalar (J) & 9.46 & \\
\emph{A ÷ I} & & \\
\bottomrule
\end{tabularx}
}

I go long another eight £5 a point blocks of S\&P 500. Because we're trading at a recent high this has the same stop level as my existing bet.

\subsection*{4 November 2014}

Sadly my long in crude is now closed.

\noindent{\scriptsize \setlength{\tabcolsep}{2pt} \renewcommand{\arraystretch}{1.12} \begin{tabularx}{\textwidth}{@{} p{0.36\textwidth}*{4}{R}@{}}
\toprule
& Crude & S\&P 500 (1) & S\&P 500 (2) & Euro Stoxx \\
\midrule
Forecast & 10 & 15 & 25 & -10 \\
Entry price & 83 & 1880 & 1980 & 2900 \\
Current price & 78.7 & 2018 & 2018 & 3034 \\
High (Low) since entry (T) & 83 & 2018 & 2018 & (2900) \\
Stop loss (S = T plus or minus R) & 79 & 1938 & 1938 & 3100 \\
Current position, blocks & 31 & 5 & 8 & -6 \\
Gain (Loss) & (£1,333) & £3,450 & £1,520 & (£804) \\
Portfolio instrument position, blocks (O) & 0.0 (Hit stop loss) & 4.81 & 8.01 & -6.43 \\
Rounded target position, blocks (P) & 0 & 5 & 8 & -6 \\
\bottomrule
\end{tabularx}
}

The net gain is £2,833, so trading capital is up to £102,833 and daily volatility target is £964. Again \textbf{price volatility} hasn't changed significantly on any instrument. Apart from closing crude there are no adjustment trades needed. I'm tempted to take profits on my only profitable trade, the S\&P 500 longs, but I can't because it isn't part of the system!

\subsection*{13 November 2014}

I decide, belatedly, to short crude at \$75 with a forecast of -10 which corresponds to a position of short 32 blocks with a stop of \$79.

\subsection*{21 November 2014}

Ouch! Euro Stoxx has also bitten the dust.

\noindent{\scriptsize \setlength{\tabcolsep}{2pt} \renewcommand{\arraystretch}{1.12} \begin{tabularx}{\textwidth}{@{} p{0.36\textwidth}*{4}{R}@{}}
\toprule
& Crude (2) & S\&P 500 (1) & S\&P 500 (2) & Euro Stoxx \\
\midrule
Forecast & -10 & 15 & 25 & -10 \\
Entry price & 75 & 1880 & 1980 & 2900 \\
Current price & 76 & 2063 & 2063 & 3130 \\
High (Low) since entry (T) & (75) & 2063 & 2063 & (2900) \\
Stop loss (S = T plus or minus R) & 79 & 1983 & 1983 & 3100 \\
Position & -32 & 5 & 8 & -6 \\
Gain (Loss) & (£320) & £4,575 & £3,320 & (£1,380) \\
Portfolio instrument position, blocks (O) & -32.8 & 4.9 & 8.2 & 0.0 Hit stop \\
Rounded target position, blocks (P) & -33 & 5 & 8 & 0 \\
\bottomrule
\end{tabularx}
}

My accumulated profit is £4,862. If I wasn't trading systematically I'd also be tempted to cut my oil short, as I don't seem to be calling crude very well.

\subsection*{28 November 2014}

My crude oil is finally paying off -- crude drops \$5 in one day! The current \textbf{price volatility} of \$1 per day for crude now looks too low, so I'm going to double it to \$2. I'm up £7,682 in accumulated profits, so the daily \textbf{volatility target} is £1,010.

\noindent{\scriptsize \setlength{\tabcolsep}{2pt} \renewcommand{\arraystretch}{1.12} \begin{tabularx}{\textwidth}{@{} p{0.36\textwidth}*{3}{R}@{}}
\toprule
& Crude (2) & S\&P 500 (1) & S\&P 500 (2) \\
\midrule
Daily volatility target (A) & £1,010 & £1,010 & £1,010 \\
Forecast & -10 & 15 & 25 \\
Entry price & 75 & 1880 & 1980 \\
Current price & 68 & 2067 & 2067 \\
Price volatility, points (F) & 2 & 20 & 20 \\
Stop loss offset (R) & 8 & 80 & 80 \\
\emph{F × X with X = 4} & & & \\
High (Low) since entry (T) & (68) & 2067 & 2067 \\
Stop loss (S) & 76 & 1987 & 1987 \\
\emph{T plus or minus R} & & & \\
Position, blocks & -32 & 5 & 8 \\
Gain / Loss £ & £2,240 & £4,675 & £3,480 \\
Instrument value volatility (I) & £20 & £100 & £100 \\
Volatility scalar (J) & 50.5 & 10.1 & 10.1 \\
\emph{A ÷ I} & & & \\
Portfolio instrument position, blocks (O) & -16.8 & 5.05 & 8.41 \\
Rounded target position, blocks (P) & -17 & 5 & 8 \\
Trade & Buy 15 & No & No \\
\bottomrule
\end{tabularx}
}

I have to reduce my crude bet and buy 15 blocks of crude to get to my target position. Primarily because of the change in \textbf{price volatility} the current position was more than 10\% away from the desired rounded \textbf{portfolio instrument position}. Note the stop loss gap in crude is also doubled to 4 × \$2 = \$8, which is added to the new low of \$68 to give a stop level of \$76, but my position is virtually halved so I have the same amount of capital at risk on this bet.

Although this looks like taking profits it is not; the system is trading automatically to keep the amount of capital at risk constant. If I hadn't adjusted the stop then the expected holding period of my bet would have shortened. Because I've adjusted the stop I also need to cut the position, or I'll have increased the risk on the trade.

This feels like a good place to stop this diary as I've shown most of the system's most interesting features. Notice how the system went against my natural instincts when I wanted to close trades too early, or let losses run. Hopefully this should show you the benefits of rigorously sticking to a position management \textbf{framework} once you've designed it.
